% Options for packages loaded elsewhere
\PassOptionsToPackage{unicode}{hyperref}
\PassOptionsToPackage{hyphens}{url}
\PassOptionsToPackage{dvipsnames,svgnames,x11names}{xcolor}
%
\documentclass[
  10pt,
]{article}
\usepackage{amsmath,amssymb}
\usepackage{iftex}
\ifPDFTeX
  \usepackage[T1]{fontenc}
  \usepackage[utf8]{inputenc}
  \usepackage{textcomp} % provide euro and other symbols
\else % if luatex or xetex
  \usepackage{unicode-math} % this also loads fontspec
  \defaultfontfeatures{Scale=MatchLowercase}
  \defaultfontfeatures[\rmfamily]{Ligatures=TeX,Scale=1}
\fi
\usepackage{lmodern}
\ifPDFTeX\else
  % xetex/luatex font selection
  % Use the default Latin Modern fonts supplied with TeX.

\fi
% Use upquote if available, for straight quotes in verbatim environments
\IfFileExists{upquote.sty}{\usepackage{upquote}}{}
\IfFileExists{microtype.sty}{% use microtype if available
  \usepackage[]{microtype}
  \UseMicrotypeSet[protrusion]{basicmath} % disable protrusion for tt fonts
}{}
\makeatletter
\@ifundefined{KOMAClassName}{% if non-KOMA class
  \IfFileExists{parskip.sty}{%
    \usepackage{parskip}
  }{% else
    \setlength{\parindent}{0pt}
    \setlength{\parskip}{6pt plus 2pt minus 1pt}}
}{% if KOMA class
  \KOMAoptions{parskip=half}}
\makeatother
\usepackage{xcolor}
\usepackage[margin=24mm]{geometry}
\setlength{\emergencystretch}{3em} % prevent overfull lines
\providecommand{\tightlist}{%
  \setlength{\itemsep}{0pt}\setlength{\parskip}{0pt}}
\setcounter{secnumdepth}{-\maxdimen} % remove section numbering
\usepackage{microtype}
\setlength{\emergencystretch}{3em}
\widowpenalty=10000
\clubpenalty=10000
\ifLuaTeX
  \usepackage{selnolig}  % disable illegal ligatures
\fi
\usepackage{bookmark}
\IfFileExists{xurl.sty}{\usepackage{xurl}}{} % add URL line breaks if available
\urlstyle{same}
\hypersetup{
  pdftitle={Rank lower bounds for node-dependent polynomial reconstruction in SAT counting},
  colorlinks=true,
  linkcolor={blue},
  filecolor={Maroon},
  citecolor={Blue},
  urlcolor={blue},
  pdfcreator={LaTeX via pandoc}}

\title{Rank lower bounds for node-dependent polynomial reconstruction in
SAT counting}
\author{Chuhang Ding (Peking University)}
\date{25 November 2026}

\begin{document}
\maketitle

\subsection{Abstract}\label{abstract}

Polynomial interpolation algorithms for counting satisfying assignments
partition the Boolean variables into groups at each interpolation point.
A factor reads its own group and linear measurements of the other
groups. We prove two rank inequalities for this interface when the
groups, measurements and factor assignments can all vary with the point.
The first bounds the total rank of measurements on a set of actual
variables by the imbalance between source and target group counts on
that set, less the degree budget of a principal determinant. A random
choice of the set gives a rank cost of order \(n/\sqrt G\) on expanding
instances when the relative degree is small, where \(G\) is the maximum
group size. The second inequality uses a tensor-valued polynomial test
and depends on how often each selected incidence crosses groups. With
balanced groups and a linear pre-substitution degree charge
\(d_0\ge c_0n\), consider the rational counting routine that visits each
Boolean-reachable measurement tuple, enumerates its prescribed Taylor
subsets and scans assignments of a largest group. The first inequality
gives a \(2^{\Omega(n^{2/3})}\) lower bound for this routine with
arbitrarily many interpolation points. The instances can be graph 2-CNFs
with at most three occurrences per variable. Under the
polynomial-dimensional schedule range of the original construction the
second inequality gives \(2^{\Omega(n\log\log n/\sqrt{\log n})}\). This
last exponent was already known for a static reconstruction template.
Here it holds for arbitrary node-dependent polynomial layouts under the
stated schedule and visit hypotheses.

\subsection{1. Introduction}\label{introduction}

Consider the task of summing a product of local constraint functions
over the Boolean cube. If each factor reads variables from just one
group, the sum factors into a product of group sums. The rational
interpolation algorithm of OpenAI {[}1{]} allows factors to cross
groups. At a nonzero interpolation point each factor reads its assigned
group and a few linear measurements from the other groups. Fixing the
vector of measurements again separates the sum. At the origin the
original factors must be recovered. A low-degree polynomial row for each
factor-variable incidence connects these requirements across points.

What cost is intrinsic to these measurements? The construction in
{[}1{]} begins with a much larger set of auxiliary coordinates. Its
placement map subsequently replaces several of those coordinates by the
same Boolean variable and fixes others to zero. Measuring all auxiliary
rows in every group can create zero rows, repeated rows, and integer
tuples that no Boolean assignment realizes. A rank bound before
placement need not survive these operations. A count of every point of a
bounding box can also exceed the size of the reachable image. We study
the matrices on the variables of a single interpolation call after all
identifications. The factor assigned to a group and the row space
available from another group may be chosen independently at every point.

\textbf{Contributions.} For one chosen output incidence per variable,
the row polynomials form a square matrix with identity constant term. We
derive two finite estimates from this matrix. A principal minor bounds
the imbalance of its input and output group counts over any set of nodes
(Theorem 1). Averaging over a random half of the variables gives a lower
bound on actual measurement rank (Corollary 2). A different argument
uses two-dimensional test spaces at the selected nodes. It yields a rank
bound in terms of each variable's number of source-target mismatches
(Theorem 4), and a rank of order \(n/(R+1)\) on some node when the
number of nodes is large compared with the square of the row degree
\(R\) (Corollary 5).

We apply these estimates to explicit operations of a counting routine.
For infinitely many bounded-occurrence graph formulas, enumerating every
Boolean-reachable measurement tuple, visiting the Taylor subsets
prescribed by a linear pre-substitution degree bound, and scanning a
largest group takes \(2^{\Omega(n^{2/3})}\) visits for all numbers of
interpolation points (Theorem 3). If the source schedule satisfies
\(r^{2Q}\le n^A\), where \(r\) is the number of groups, \(Q\) the number
of points and \(A\) a fixed constant, the bound strengthens to
\(2^{\Omega(n\log\log n/\sqrt{\log n})}\) (Theorem 6). A parsimonious
extension carries an arbitrary input formula into an instance with the
required expanding scopes.

\textbf{Previous barriers.} Dias and Lima {[}2, Lemma 6.2 and Theorems
7.1 and 8.1{]} use a product of the own coefficients of incidences to
constrain the structure of a formula at an interpolation point. Among
their bounds is \(2^{\Omega(\sqrt{\operatorname{tw}})}\) for the
clause-product scheme in a model that also treats interaction-graph
dynamic programming. Our principal determinants include all entries of a
selected square submatrix. Their root multiplicities control rank
deficiencies, which can be bounded using the imbalance of source and
target counts.

Tulsyan {[}3, Theorem 5.1 and Proposition 5.2{]} proves a rank lower
bound using a triangular submatrix of \(I-M_0\), where a single matrix
\(M_0\) satisfies a Hamming-distance support constraint. With the
Arora--Rao--Vazirani structure theorem, his result gives the exponent
\(n\log\log n/\sqrt{\log n}\) on linked bounded-occurrence formulas
across arbitrary schedules in that template. His companion {[}4,
Sections 5--7{]} distinguishes the static template from the more
permissive point-dependent reconstruction condition. Our
restricted-schedule theorem attains this exponent for the latter
condition on physical variable groups. Its numerical schedule
restriction remains essential to the proof. The arbitrary-schedule
theorem gives the exponent \(n^{2/3}\) under explicit visits.

The main theorem of {[}1{]} uses a separate tensor contraction in
Appendix D. Its computation does not enumerate the Taylor subsets or
rational measurement images considered here. All runtime bounds below
concern the stated interpolation routine.

\subsection{2. Actual-variable layouts and the work
measure}\label{actual-variable-layouts-and-the-work-measure}

Let there be \(N\) Boolean variables and a finite list of factors with
specified scopes. Let \(\lambda_1,\ldots,\lambda_Q\) be distinct nonzero
rational numbers. At point \(t\) let \(S_{t,1},\ldots,S_{t,r}\)
partition the variables, allowing empty groups. Write \(a_v(t)\) for the
group of variable \(v\), and \(b_f(t)\) for the group assigned to factor
\(f\). Group \(g\) has a rational measurement matrix \(C_{t,g}\) with
columns indexed by \(S_{t,g}\). An incidence \((f,v)\) has a polynomial
row \(P_{f,v}(X)\) of degree at most \(R\) such that

\[
 P_{f,v}(0)=e_v^\top,\qquad
 P_{f,v}(\lambda_t)|_{S_{t,g}}\in\operatorname{rowspan}C_{t,g}
             \quad\text{whenever } g\ne b_f(t).                       \tag{1}
\]

This is the linear part of Definition 3.1 of {[}1{]}. The source
additionally bounds coefficient heights for bit complexity. The two rank
inequalities allow arbitrary rational coefficients. Their hypotheses
impose no relation between the measurement matrices at different nodes.

Fix a nonempty set \(V\) of variables and an incident factor \(f(v)\)
for each \(v\in V\). With rows \(P_{f(v),v}\) and input columns \(V\),
form \(P(X)\in\mathbb Q[X]^{V\times V}\). Write \(b_v(t)=b_{f(v)}(t)\).
For \(U\subseteq V\) let

\[
 m_t(U)=\sum_{g=1}^{r}\operatorname{rank}\!\left(
                  C_{t,g}|_{U\cap S_{t,g}}\right).
\]

The restriction always means retaining all rows and the indicated
columns. With \(C_t\) the block-diagonal matrix of the group maps
restricted to \(V\), (1) gives

\[
 P(0)=I_V,\qquad P(\lambda_t)=D_t C_t+L_t,\qquad
 L_t[u,v]=0\quad\text{if }a_v(t)\ne b_u(t).                          \tag{2}
\]

To see this, represent the restricted coefficients of each output on
each foreign group as a linear combination of that group's measurement
rows. Subtracting these combinations leaves an output supported on its
assigned group. Further restriction of inputs and outputs to \(U\)
preserves the decomposition.

The jointly reachable measurement image at \(t\) is

\[
  \mathcal Y_t=\{(C_{t,g}x|_{S_{t,g}})_{g=1}^{r}:x\in\{0,1\}^N\}.
\]

For each \(U\subseteq V\), its cardinality is at least \(2^{m_t(U)}\).
Indeed, choose independent columns of maximal number in each restricted
group matrix. Their Boolean subset sums are distinct, and choices in
disjoint groups multiply. One may fix all other bits to zero.

We state the cost assumptions separately from the algebraic model. Let
\(T_{\rm run}\) count elementary record visits, or operations that
charge at least a constant for each visit. In the balanced model
\(|S_{t,g}|\le2N/r\). The three operations charged are as follows.

\begin{enumerate}
\def\labelenumi{\arabic{enumi}.}
\tightlist
\item
  At every node the routine visits each member of \(\mathcal Y_t\) at
  least once. A tuple is reachable if some Boolean input produces it,
  even when the final weighted sum on its fiber is zero.
\item
  At some node the routine visits each subset of \([N]\) of size at most
  \(s=\min(N,\lfloor d_0R/Q\rfloor)\), where \(d_0\ge c_0N\) for fixed
  \(c_0>0\).
\item
  At some partition it scans all Boolean assignments on a largest group.
\end{enumerate}

The first operation can cost substantially less than enumerating the
full integer box used in the source. The second is the subset visit of a
fixed-substitution Taylor evaluator. For the product of separate binary
clause indicators, \(d_0\) can be taken as the product degree in
independent incidence indeterminates. We assume the indicated
enumeration even if the final univariate polynomial for a particular
input happens to have smaller degree. The third operation costs at least
\(2^{N/r}\) visits. The nodes witnessing the three costs may differ.

\subsection{3. Principal determinants and group
imbalance}\label{principal-determinants-and-group-imbalance}

For fixed outputs \(f(v)\), let \(A_g^U(t)\) count \(v\in U\) with
source \(a_v(t)=g\), and let \(B_g^U(t)\) count \(v\in U\) with target
\(b_v(t)=g\).

\textbf{Theorem 1 (principal-subset imbalance).} For every
\(U\subseteq V\) and every set \(T\) of nodes,

\[
 \sum_{t\in T}m_t(U)+R|U|
 \ \ge\ \frac12\sum_{t\in T}\sum_g|A_g^U(t)-B_g^U(t)|.             \tag{3}
\]

\textbf{Proof.} The principal submatrix \(P_{U,U}(X)\) has identity
constant term. Its determinant is nonzero and has degree at most
\(R|U|\). At node \(\lambda_t\) its order of vanishing is at least the
nullity of \(P_{U,U}(\lambda_t)\). To verify this, apply invertible
constant row operations that make as many evaluated rows zero as the
nullity. Every such polynomial row contributes a factor \(X-\lambda_t\)
to the determinant. Distinct nodes contribute relatively prime factors,
giving

\[
 \sum_{t\in T}\operatorname{nullity}P_{U,U}(\lambda_t)\le R|U|.      \tag{4}
\]

After grouping the rows of \(L_{t;U,U}\) by target and the columns by
source, the remaining blocks at group \(g\) are \(B_g^U(t)\) by
\(A_g^U(t)\). Its rank is at most \(\sum_g\min(A_g^U(t),B_g^U(t))\).
Both group-count sequences sum to \(|U|\), so this rank bound equals
\(|U|-\tfrac12\sum_g|A_g^U(t)-B_g^U(t)|\). The measured part of (2) on
\(U\) has rank at most \(m_t(U)\). Adding the two rank bounds,
rearranging and using (4) proves (3). The empty principal set gives zero
on both sides. \(\square\)

For a graph application, we choose \(f(v)\) using a bounded set of
candidate incidences. Say that \(V\), of size \(n\), has a collection of
candidate factors with parameters \(\delta,\Delta\) if each \(v\) has a
nonempty list of at most \(\Delta\) incident factors and at every node
at least \(\delta n\) pairs on these lists have \(a_v(t)\ne b_f(t)\).
Write \(\beta=\delta/\Delta\).

\textbf{Corollary 2 (rank after restricting a principal set).} If no
group contains more than \(G\) variables of \(V\), then

\[
 \frac1Q\sum_t m_t(V)+\frac{nR}{2Q}
    \ \ge\ \frac{\beta n}{48\sqrt{2G}} .                            \tag{5}
\]

\textbf{Proof.} For each variable, select the factor on its list that
maximizes its number of mismatches over all nodes. Comparing this
maximum with its list average gives at least \(\beta nQ\) mismatches,
summed over selected variables and nodes. Now choose \(U\) by keeping
each variable independently with probability \(1/2\).

Fix one node and group. Among the selected incidences on \(V\), let
\(x\) count those starting in that group and ending elsewhere, and \(y\)
those ending there and starting elsewhere. The two collections involve
disjoint variables. Matched incidences cancel from the group-count
difference. Consequently

\[
 A_g^U(t)-B_g^U(t)\ \stackrel{\rm law}{=}\ (x-y+S_{x+y})/2 ,
\]

where \(S_h\) is the sum of \(h\) independent fair signs. Symmetry and
the triangle inequality give \(\mathbb E|z+S_h|\ge\mathbb E|S_h|\) for
every \(z\). For \(h\ge1\), \(\mathbb E S_h^2=h\) and
\(\mathbb E S_h^4=3h^2-2h\le3h^2\). On the event \(E=\{S_h^2\ge h/2\}\)
one has \(\mathbb E[S_h^2{\bf1}_E]\ge h/2\). Cauchy--Schwarz gives
\(\Pr(E)\ge1/12\). It follows that

\[
 \mathbb E|A_g^U(t)-B_g^U(t)|
     \ge\frac{\sqrt{x+y}}{24\sqrt2}
     \ge\frac{x}{24\sqrt{2G}} ,
\]

also when \(x+y=0\). The last step uses \(x\le G\). Summing over all
groups counts each mismatch once in \(\sum x\). Apply (3), take
expectations, and use \(\mathbb E|U|=n/2\) and \(m_t(U)\le m_t(V)\).
This yields (5). \(\square\)

\textbf{Theorem 3 (work for arbitrary node counts).} There is a sequence
of positive 2-CNF formulas on \(n\) variables, each variable occurring
at most three times, for which any balanced layout satisfying the three
visit assumptions has

\[
                        T_{\rm run}\ge 2^{\Omega_{c_0}(n^{2/3})}.       \tag{6}
\]

The number of interpolation points, the row heights and the auxiliary
dimension are unrestricted.

\textbf{Proof.} Appendix A gives a family of simple graphs of degrees
two or three such that every cut with smaller side size between \(n/4\)
and \(n/2\) has at least \(\kappa n\) edges, for fixed \(\kappa>0\).
Introduce the factor \(x_u\lor x_v\) for every graph edge \(uv\). If
\(r<4\), a group scan costs \(2^{\Omega(n)}\). Otherwise each group has
at most \(n/2\) variables. At each node some union of groups has size in
\([n/4,n/2]\). If no individual group has such a size, accumulate groups
of size below \(n/4\) until the lower endpoint is reached. The endpoints
of each edge crossing this cut belong to different groups. No matter
where its clause is assigned, at least one endpoint mismatches its
target. Using all incident graph edges as lists gives
\(\delta=\kappa,\Delta=3\).

Put \(G=2n/r\), and fix \(c_*=1/(48\sqrt2)\). If \(G\ge n^{2/3}\), any
largest group has at least \(G/2\) variables, which proves (6). Suppose
\(G<n^{2/3}\). If \(R/Q\le c_*\beta/\sqrt G\), (5) gives a node with
\(m_t(V)\ge c_*\beta n/(2\sqrt G)=\Omega(n^{2/3})\). The visit to its
reachable image proves (6). Finally, if \(R/Q>c_*\beta/\sqrt G\), then
\(d_0R/Q=\Omega_{c_0}(n^{2/3})\). The prescribed Taylor cutoff contains
\(s=\Omega_{c_0}(n^{2/3})\) variables even after the cap at \(n\) and
rounding. All subsets of any \(s\) variables are visited, giving at
least \(2^s\) visits. \(\square\)

For these formulas the pre-substitution degree in separate incidence
variables is exactly \(2|E(G)|\ge2n\). The quadratic term of each clause
indicator has nonzero coefficient and involves two variables that belong
to no other factor. The layout argument also holds for arbitrary literal
signs on the same scopes. The positive formulas alone have a trivial
affirmative SAT answer. Theorem 3 concerns their counting routine under
the stated visits.

\subsection{4. A polynomial test on selected
nodes}\label{a-polynomial-test-on-selected-nodes}

For a set \(T\) of \(k\) nodes and a chosen output incidence at \(v\),
write \(h_v=\#\{t\in T:a_v(t)\ne b_v(t)\}\). Define

\[
 \phi_R(h)=2^{-h}\sum_{i=0}^h\binom hi(h-R-2i)_+,\qquad
 (z)_+=\max(z,0).
\]

\textbf{Theorem 4 (mismatch rank bound).} Whenever \(k>R\),

\[
                 \sum_{t\in T}m_t(V)\ \ge\ 2\sum_{v\in V}\phi_R(h_v).
                                                                         \tag{7}
\]

The chosen incidence per variable may depend on \(T\). At \(k=1,R=0\),
the right side is exactly the number of variables whose selected
incidence is foreign. Their constant rows demand all of those
coordinates from their source groups.

\textbf{Proof.} Let \(E=\mathbb Q^2\) and \(W=E^{\otimes k}\), with one
factor per \(t\in T\). Write \(M=2^k\), and choose distinct rational
lines \(\ell_g\subset E\) indexed by group labels. Let \(\ell^\circ\)
denote the line of covectors annihilating \(\ell\). Each physical
variable \(v\) is given two bases, one of \(W\) and one of \(W^*\), by
making these choices at every \(t\).

\begin{itemize}
\tightlist
\item
  When \(a=a_v(t)\ne b=b_v(t)\), take input low and high lines
  \(\ell_a,\ell_b\) and output low and high lines
  \(\ell_b^\circ,\ell_a^\circ\). A local input pairs only with the
  output of the same index.
\item
  When \(a=b\), choose a complement \(\ell'\) to \(\ell_a\). Take input
  lines \(\ell_a,\ell'\) and output lines
  \(\ell_a^\circ,(\ell')^\circ\). In this case it is the opposite
  indices that pair.
\end{itemize}

Choose basis vectors and scale covectors so the nonzero pairings equal
one. In the tensor bases let \(e_{v,S}\) and \(f_{v,B}\) denote the
elements with high sites \(S\) and \(B\). The unique output dual to
\(e_{v,S}\) has index

\[
 B_v(S)=S\mathbin{\triangle}\{t\in T:a_v(t)=b_v(t)\}.                \tag{8}
\]

The bases can depend on \(v\) because every occurrence of a physical
operator \(P,C_t,L_t\) below is its tensor product with the identity on
the same ambient \(W\).

Write \(\pi_S(X)=\prod_{t\in S}(X-\lambda_t)\). For \(j\ge0\), let
\(H_j\) consist of tuples of \(W\)-valued polynomials \(h_v(X)\) of
degree at most \(j\), whose scalar coordinate along \(e_{v,S}\) is
divisible by \(\pi_S\). Define

\[
 I_j=\{h\in H_j:(C_t\otimes I_W)h(\lambda_t)=0\text{ for each }t\in T\},
 \qquad I_{-1}=H_{-1}=0 .
\]

At the evaluation point \(\lambda_t\), \(h_v(\lambda_t)\) only has terms
whose \(t\)-factor is on the input low line \(\ell_{a_v(t)}\). For all
variables in group \(g\) this is a common subspace \(W_{t,g}\) of
dimension \(M/2\). The contribution of that group to the evaluated
measurement factors through
\(C_{t,g}|_{V\cap S_{t,g}}\otimes I_{W_{t,g}}\), of rank
\((M/2)\operatorname{rank} C_{t,g}|_{V\cap S_{t,g}}\). Stacking these
constraints takes at most the sum of their ranks over all groups and
nodes. Hence

\[
               \dim H_j-\dim I_j\le (M/2)\sum_{t\in T}m_t(V).          \tag{9}
\]

We next obtain a lower bound on this codimension. Let \(J=k-R-1\),
\(0\le j\le J\), and \(h\in I_j\). For each \(v,B\), form the scalar
polynomial

\[
      \pi_B(X)f_{v,B}\!\left(((P(X)\otimes I_W)h(X))_v\right).
\]

It vanishes at every \(\lambda_t\). If \(t\in B\) this follows from
\(\pi_B\). If \(t\notin B\), the measured part of (2) vanishes by
\(h\in I_j\). In the remaining terms \(L_t[v,w]\ne0\) requires
\(a_w(t)=b_v(t)\). Their input factor at \(t\) therefore lies on the
line annihilated by the output low covector. When \(|B|\le J-j\), the
scalar polynomial has degree at most \(|B|+R+j<k\) and is identically
zero. Evaluating at the origin and using \(P(0)=I\) and \(\pi_B(0)\ne0\)
shows that the \((v,S)\) coordinate of \(h(0)\) vanishes whenever
\(|B_v(S)|\le J-j\). Divisibility also eliminates all coordinates with
\(|S|>j\).

Evaluation at zero on \(I_j\) has kernel \(XI_{j-1}\). Dividing by \(X\)
preserves the divisibility by each \(\pi_S\) and the nodal kernel
constraints because every interpolation point is nonzero. Thus the
increment \(\dim I_j-\dim I_{j-1}\) is at most

\[
 \#\{(v,S):|S|\le j\}
 -\#\{(v,S):|S|\le j,\ |B_v(S)|\le J-j\}.
\]

For \(H_j\) its dimension increment is exactly the first count. Each
scalar coordinate along \(e_{v,S}\) contributes \((j-|S|+1)_+\)
dimensions to \(H_j\). Summing the increment inequality through \(J\)
and using (9) yields

\[
 (M/2)\sum_{t\in T}m_t(V)
      \ge \sum_{v,S}(k-R-|S|-|B_v(S)|)_+ .
\]

For a uniform \(S\), the \(k-h_v\) matched sites each contribute one to
\(|S|+|B_v(S)|\) by (8). The remaining sites contribute twice a
\(\operatorname{Bin}(h_v,1/2)\) random variable. Divide by \(M/2\) to
obtain (7). \(\square\)

\textbf{Corollary 5 (many nodes with large rank).} Under the
candidate-factor hypothesis with \(\beta>0\), there are constants
\(c_\beta,C_\beta>0\) such that for
\(k=\lceil C_\beta(R+1)^2\rceil\le Q\), every set \(T\) of \(k\) nodes
satisfies

\[
                    \frac1k\sum_{t\in T}m_t(V)
                      \ge c_\beta\frac{n}{R+1}.                       \tag{10}
\]

At least \(Q-k+1\) nodes meet this rank lower bound individually.

\textbf{Proof.} Select \(f(v)\) to maximize mismatches on this
particular \(T\), giving \(\sum_v h_v\ge\beta nk\). At least
\(\beta n/2\) variables have \(h_v\ge\beta k/2\), since otherwise their
total falls below this bound even with \(h_v\le k\). With
\(C_\beta=16/\beta\), those \(h_v\) satisfy \(h_v\ge8(R+1)^2\). The
fourth-moment computation in Corollary 2 and symmetry give
\(\Pr(S_h\ge\sqrt{h/2})\ge1/24\) for \(h\ge1\). Thus

\[
        \phi_R(h)=\mathbb E(S_h-R)_+\ge \sqrt h/(48\sqrt2)
         \quad\text{when }h\ge8(R+1)^2.
\]

Equation (7) now bounds \(\sum_{t\in T}m_t(V)\) below by a constant
depending only on \(\beta\) times \(n\sqrt{k}\). Since
\(k\le(C_\beta+1)(R+1)^2\), this proves (10). If \(k\) nodes strictly
violated the individual bound, their average would contradict (10).
\(\square\)

\subsection{5. The polynomial-dimensional
schedule}\label{the-polynomial-dimensional-schedule}

The source's full auxiliary schedule has \(r^{2Q}\) positions. Its
polynomial-dimensional parameter choice imposes a numerical condition on
\(r,Q\). The next result needs the condition only; all ranks in its
proof remain on physical variables.

\textbf{Theorem 6.} On the graph formula family of Theorem 3, assume a
balanced layout, so that \(|S_{t,g}|\le 2n/r\) at every node \(t\) and
group \(g\), and the same three explicit visit assumptions. If
\(r^{2Q}\le n^A\) for fixed \(A>0\), then

\[
       \log T_{\rm run}=\Omega_{A,c_0}\!\left(
                       n\frac{\log\log n}{\sqrt{\log n}}\right).        \tag{11}
\]

\textbf{Proof.} Use natural logarithms and put \(L=\ln n,\ l=\ln L\).
For \(r\le\sqrt L/l\), the group scan costs at least \(2^{n/r}\), giving
(11). Otherwise, for large \(n\), \(r\ge4,\ \ln r\ge l/3\), and
\(Q\le(3A/2)L/l\). The balanced-cut argument supplies a fixed
\(\beta>0\).

For sufficiently large \(Q\), set \(u=\sqrt{Q/\ln(eQ)}\). If
\(R+1\le u\), then \(\lceil C_\beta(R+1)^2\rceil\le Q\) for all
sufficiently large \(Q\). By Corollary 5 and the reachable-image bound,
\(\ln T_{\rm run}\ge c n/(R+1)\ge c n/u\) for a fixed \(c>0\). If
\(R+1>u\), then \(R\ge u/2\). Fix \(a>0\) sufficiently small in terms of
\(c_0\) and put \(s_0=\lfloor an/\sqrt{Q\ln(eQ)}\rfloor\). The Taylor
cutoff is at least \(s_0\). Also \(s_0\le n/2\) and \(s_0\) tends to
infinity uniformly for \(Q=O(\ln n)\). Therefore

\[
       \ln T_{\rm run}\ge\ln\binom n{s_0}
        \ge s_0\ln(n/s_0)=\Omega\!\left(n\sqrt{\ln(eQ)/Q}\right).
\]

For \(Q\) below the fixed threshold and \(R\ge1\), the Taylor subsets
already give a linear exponent. For \(R=0\) all selected incidence rows
are constant coordinate rows. Every variable with some mismatched
incidence must have its coordinate row in the measurement row space of
its source group. There are at least \(\delta n/\Delta\) such variables
at each node, giving a linear rank and hence a linear exponent. Finally
\(x\mapsto\sqrt{\ln(ex)/x}\) decreases on \(x\ge1\). Substituting the
upper bound \((3A/2)L/l\) proves (11). \(\square\)

\subsection{6. Scope and verification}\label{scope-and-verification}

The rank inequalities (3) and (7) concern linear recovery on the actual
rational variable space. They are unaffected by deleting dependent
measurement rows or by the construction procedure for the final layout.
The time corollaries need the three operations of Section 2. A procedure
can avoid the charged work by aggregating many measurement images in one
operation, by using additional knowledge that particular fibers have
zero contribution, or by recovering Taylor data without the prescribed
subset enumeration. An unrestricted solver can also avoid the
interpolation call entirely. Theorems 3 and 6 do not address those
procedures or imply a general SAT lower bound.

For Theorem 3 the balance condition can be dispensed with if a largest
group is fully scanned at every node. If a group ever contains
\(n^{2/3}\) variables, its scan suffices. Otherwise every group is
smaller than \(n^{2/3}\le n/2\) for large \(n\), so expansion works at
every node. Applying (5) with \(G=n^{2/3}\) gives the same
rank-or-Taylor dichotomy.

The unrestricted-node analogue of (11) remains open here. Another
question is whether a lower bound on the measured row spaces can be
converted into a lower bound on joint evaluation when a single operation
can sum many fibers. Such a statement requires a cost model specifying
the joint operation.

The supplemental checker constructs feasible polynomial matrices first
and derives the minimal measurement row spaces containing their
foreign-group restrictions. Using exact integer elimination, it verified
61,824 mismatch inequalities on node and variable subsets and 120,864
principal-minor and imbalance comparisons. It also checked 38 explicit
tensor filtration levels, 44,352 restricted Boolean images, and 1,096
unique gate extensions. The 160 sampled layouts have up to five
variables and six nodes. The accompanying Lean development formalizes the principal-subset and
mismatch rank inequalities, their probabilistic corollaries, the graph
family, and the visit bounds in Theorems 3 and 6. The repository
provides the theorem index, the pinned build and the axiom reports.
The finite layout checker is included as
\texttt{tests/actual\_layout\_check.py}.

\textbf{Contributions and review status.} Chuhang Ding set the research questions and direction and proposed conjectures. AI assistants and agents contributed to proofs, tests, Lean formalization, review, and writing. This is an unsubmitted working manuscript pending
independent human mathematical review.

\subsection{Appendix A. Bounded-degree graphs and a parsimonious
extension}\label{appendix-a.-bounded-degree-graphs-and-a-parsimonious-extension}

\textbf{Lemma A.1.} There is a fixed \(\kappa>0\) and, for every
positive \(n\) divisible by \(128\), a simple graph on \(n\) vertices of
degrees two or three for which each cut with smaller side in
\([n/4,n/2]\) has at least \(\kappa n\) crossing edges.

\textbf{Proof.} Put \(m=n/32\), so \(4\mid m\). Choose sixteen
independent uniform permutations of \(m\) vertices. For \(B\) of size
\(m/4\le b\le m/2\), the probability that a permutation sends at most
\(b/8\) elements of \(B\) outside \(B\) is at most

\[
 \binom b{\lfloor b/8\rfloor}(b/m)^{\lceil7b/8\rceil}
       \le 2^{-(7/8-H_2(1/8))b},
\]

where \(H_2\) is binary entropy. Some subset of \(\lceil7b/8\rceil\)
specified vertices would all have to land in \(B\). Each candidate
subset has probability at most \((b/m)^{\lceil7b/8\rceil}\) by sampling
without replacement. Since \(7/8-H_2(1/8)>1/4\), taking the sixteenth
power and summing over at most \(2^m\) choices gives total probability
less than one. Thus for every such \(B\) one permutation sends more than
\(b/8\) elements outside. Convert all permutation arcs to undirected
edges and discard loops and duplicates. The outward arcs of the
particular permutation have distinct inside endpoints, so represent
distinct crossing edges. The graph has maximum degree \(32\) and at
least \(m/32\) edges crossing each required cut.

Replace each old vertex by a cycle of length \(D=32\), attaching each
old incident edge at a different vertex on the cycle. The resulting
simple graph has \(Dm=n\) vertices of degrees two or three. For a side
\(\mathcal U\) of size in \([n/4,n/2]\), denote by \(p\) the number of
cycles split by it. If \(p\ge m/(64D)\), cycle edges supply at least
\(2p\ge m/(32D)\) crossings. Otherwise round the side to an old vertex
subset \(B\) containing the full cycles and contained in the full plus
split cycles, of size between \(m/4\) and \(m/2\). Such a count exists
because \(|\mathcal U|/D\) lies between the two integral cycle counts
and in the interval with integral endpoints. At most \(Dp\) incident old
edges can change their crossing status in the rounding. The original cut
in the new graph therefore has at least \(m/32-Dp>m/64\) crossings. Both
cases allow \(\kappa=1/(32D^2)\). \(\square\)

\textbf{Proposition A.2.} The rank and visit conclusions can be imposed
on a parsimonious extension of any CNF, with a linear number of added
variables and clauses.

\textbf{Proof.} For a formula \(F\) on a supported number \(n\) of
original variables \(V\), choose the graph of Lemma A.1 and add a
variable \(y_{uv}\) for each edge \(uv\). With \(y=y_{uv}\), add

\[
      (\neg y\lor x_u),\qquad(\neg y\lor x_v),\qquad
      (y\lor\neg x_u\lor\neg x_v).
\]

Each original assignment has the unique extension \(y_{uv}=x_ux_v\). The
extended formula has exactly the same number of satisfying assignments.
It has \(N=n+|E|\le5n/2\) variables and \(O(n)\) additional clauses.
Each old variable gains at most six occurrences and each new variable
has three. Width increases to at most
\(\max(3,\operatorname{width}(F))\). Other original sizes can be padded
to a multiple of \(128\) using at most \(127\) extra variables pinned by
unit clauses. The construction takes linear time when the graph is
supplied. Lemma A.1 itself is an existence proof.

When \(r\ge10\), a group contains at most \(2N/r\le n/2\) original
variables. Find a balanced union of groups on \(V\) alone, and use the
ternary clause of each graph edge as the candidate factor for its two
original variables. Every crossing edge gives a mismatch at an original
endpoint, and an original variable has at most three such candidate
factors. Thus \(V\) satisfies the same fixed list parameters. Use
\(G=2N/r\), with \(n=\Theta(N)\), in the proof of Theorem 3 and use
Corollary 5 on \(V\) in Theorem 6. Ranks restricted to \(V\) bound the
reachable image from below even with all fresh bits fixed to zero. If
\(r<10\), the group scan already has a linear exponent. The runtime
conclusions follow in the total variable count \(N\) under their
respective visit and schedule assumptions. A solver that removes the
defining gate clauses by preprocessing need not encounter this layout.
\(\square\)

\subsection{References}\label{references}

\begin{enumerate}
\def\labelenumi{\arabic{enumi}.}
\tightlist
\item
  OpenAI. \emph{The Exponential Time Hypothesis is False}. Revised
  October 25, 2026. Sections 3--4 describe the rational counting routine
  and the incidence layout. Appendix D gives the tensor-contraction
  algorithm.
  \href{https://github.com/openai/math/tree/main/preprints/The-Exponential-Time-Hypothesis-is-False-October-25-2026}{Manuscript
  and source}
\item
  Fabio Dias and Yvens Lima. \emph{Limits and Constants of OpenAI's
  Claimed ETH Refutation}. arXiv:2611.00517v1, 2026.
  \href{https://arxiv.org/html/2611.00517v1}{Full text}
\item
  Vaibhav Tulsyan. \emph{What the ETH Refutations Do Not Give: Barriers
  for the Counting Engines, with Machine-Checked Lemmas}. Hexagon
  2611.00149v1, 2026.
  \href{https://hexagonmath.org/pdf/2611.00149v1}{Paper}
\item
  Vaibhav Tulsyan. \emph{The parity barrier in OpenAI's Engine I: exact
  constraints, a lower bound for every \(M_0\), and what is left open}.
  Companion research note for {[}3{]}, Sections 5--7, repository
  revision 946c49993b78aa163fcac3f8cb78ce351dd5dcc6.
  \href{https://github.com/xennygrimmato/eth-hamiltonicity/blob/946c49993b78aa163fcac3f8cb78ce351dd5dcc6/research/permanent/parity-barrier.md}{Pinned
  source}
\end{enumerate}

\end{document}
